\documentclass[11pt]{article}
% === core_packages.tex ===
\usepackage{amsmath, amssymb, amsthm, amsfonts}
\usepackage{geometry}
\usepackage{hyperref}
\usepackage{bookmark}
\usepackage{graphicx}
\usepackage{physics}
\usepackage{booktabs}
\usepackage{listings}
\usepackage{amscd}
\usepackage{tikz-cd}
\usepackage{mathtools}
\usepackage{xspace}
\usepackage[backend=biber, style=numeric]{biblatex}
\usepackage{tcolorbox}
\tcbuselibrary{theorems}
\usepackage{parskip}
\usepackage{tikz}
\usepackage{etoolbox}
\usetikzlibrary{arrows.meta, shapes, positioning, calc, cd, shapes.geometric, backgrounds}
\usepackage[en-US]{datetime2}
\usepackage{array}

% --- Theory names ---
\newcommand{\iame}{[theory needs name]}
\newcommand{\iameverything}{[theory needs name]}
\newcommand{\theory}{[theory needs name]}


% --- Math & Ontology Macros ---
\newcommand{\entropy}{S}
\newcommand{\opponent}{opponent}
\newcommand{\aside}[1]{[\emph{\opponent}: ``#1'']}
\newcommand{\paper}[1]{[\emph{#1}: ``#1'']}
\newcommand{\fabric}{voxel}
\newcommand{\Fabric}{Voxel}
\newcommand{\pdflink}[1]{}



% --- Global Layout Defaults ---
\geometry{letterpaper, margin=1in}
\hypersetup{
    colorlinks=true,
    linkcolor=black,
    citecolor=black,
    filecolor=magenta,
    urlcolor=blue
}
\hfuzz=5pt \vfuzz=5pt \hbadness=10000 \vbadness=10000

\author{\iame}


\usepackage{tocloft}

% --- Article Specific Theorems ---
\newtheorem{lemma}{Lemma}
\theoremstyle{definition}
\newtheorem{axiom}{Axiom}
\newtheorem{theorem}{Theorem}
\newtheorem{corollary}{Corollary}
\newtheorem{definition}{Definition}
\newtheorem{proposition}{Proposition}
\newtheorem{conjecture}{Conjecture}

\makeatletter
\let\remark\@undefined
\let\endremark\@undefined
\makeatother
\theoremstyle{remark}
\newtheorem*{remark}{Remark}

\newtcbtheorem[number within=section]{principle}{Principle}%
{colback=blue!5,colframe=blue!50!black,fonttitle=\bfseries}{principle}

% --- Article TOC Layout ---
\setlength{\cftbeforesecskip}{1.0em}
\renewcommand{\cftsecleader}{\cftdotfill{\cftdotsep}}
\renewcommand{\cftsubsecleader}{\cftdotfill{\cftdotsep}}



% Base URL of the folder that holds the Python tests 
\newcommand{\repo}{https://github.com/juhameskanen/abstract/blob/main/tests}
% \testfile{name.py} typesets the file name as a link into \repo
\newcommand{\testfile}[1]{\href{\repo/#1}{\texttt{\detokenize{#1}}}}


\newcommand{\flag}[1]{\textbf{[#1]}}

\newtheorem{assumption}{Assumption}


\addbibresource{../references.bib}

\title{Spectral Compression and the Emergence of Gravity:\\
A Consistency Study of the Heat-Kernel Structure of a Spectral Description Cost}
\author{Juha\\[2pt]\small The Abstract Universe Project}
\date{Draft, 7 November 2026}

\begin{document}

\maketitle




\begin{abstract}
The Wavefunction Compression Principle asserts that the quantum wavefunction is the
universe's data-compression scheme, with the probability of a state falling as $2^{-C_s}$,
where $C_s$ is a \emph{spectral complexity} (a description cost counted over spectral modes).
This study asks whether such a cost, once the geometry is promoted to a variable, contains
general relativity. We study the functional $C[g]=\Tr f(L_g/\Lambda^2)$ for a smooth cutoff
$f$ and a Dirac-type operator $L_g$, and test its heat-kernel structure numerically against
exact or highly accurate spectra on spheres, products of spheres, tori, and a planar
inhomogeneous background. Under the explicitly listed assumptions, we find: (i) the volume and
Einstein--Hilbert terms appear with the standard coefficients, and the sign of the
Einstein--Hilbert term is correct only for spinor-type operators; (ii) a magnetic flux
background acts as a source whose effect on the stationary geometry matches Einstein--Maxwell,
with $G/e^2$ fixed by the cutoff moments; (iii) the Maxwell term follows from varying the gauge
field and is proportional to $f(0)$; (iv) with the per-mode linear cost applied only to occupied modes and
the identification $\beta=\kappa/\Delta\omega$, the stationary conditions on $S^3\times S^1$ reproduce
the Einstein static universe; (v) in a planar weak-field setting the response reproduces
$\nabla^2\Phi=4\pi G\rho$ with the same $G$. These are \emph{consistency checks of the local
spectral action with a specific matter coupling}, not a derivation of gravity from the
compression principle: several inputs (cutoff shape, a hand-tuned cosmological counterterm,
Euclidean signature, the occupied-mode prescription, $D=4$) remain assumed. All failures and
corrections encountered along the way are logged in Section~\ref{sec:failures}.
Every numerical claim is tied to a script in the accompanying repository.
\end{abstract}



\section{Starting point and scope}\label{sec:start}

\subsection{The compression paper}
The starting text (chapter ``The Deep Nature of Quantum Mechanics'') proposes that the
wavefunction is a dense spectral encoding, analogous to a DCT in video compression, and replaces
Kolmogorov complexity by a computable \emph{spectral complexity}
\begin{equation}
C_s(\Psi)=\mathrm{Cost(base)}+\sum_{i=1}^{N}\Big[\mathrm{Cost}(\phi_i)+\mathrm{Cost}(A_i)
+\frac{\omega_i}{\Delta\omega}\Big],
\label{eq:Cs}
\end{equation}
with Solomonoff \cite{solomonoff1964} weights $P\propto 2^{-C_s}$, i.e.\ a Boltzmann factor in which $\Delta\omega$
plays the role of a temperature. Simulations (not reproduced here) starting from a minimal-entropy
state, with spectral cost and an observer filter, produce smooth evolving structure,
interference patterns and lattice-like symmetric structures. The paper leaves open the origin of
gravitational curvature.

\subsection{Inputs assumed, but not tested here}
The following are statements either assumed, or derived earlier of the framework.
Nothing in the present note derives or tests them.
\begin{itemize}
\item Spacetime is $D=4$ (this is assumed for now; believed to emerge as the most probable dimension).
\item The counted object is the whole 4D history (spectral cost over the entire walk).
\item The probability of emergent microstructure is lognormal in time, giving a three-fold
expansion curve; we sit near the peak where probability of emergence of matter is highest, and the Hubble rate is therefore minimal.
\end{itemize}

\subsection{Scope}
We ask a narrow question: \emph{if the spectral cost is a functional of the geometry, what does
it contain?} We do not address the observer filter, the origin of the cutoff, Lorentzian
signature, or the selection of dimension.



\section{Framework}
\label{sec:framework}

\subsection{Spectral cost of a geometry}

Let $L_g=-(\nabla^2+E)$ be a Laplace-type operator on a closed Riemannian 4-manifold $(M,g)$,
and $f$ a smooth, rapidly decaying cutoff profile, $\Lambda$ the cutoff scale ($\Lambda=1$ in
all numerics). Define

\begin{equation}
C[g]=\Tr f(L_g/\Lambda^2).
\end{equation}

For smooth $f$ the standard heat-kernel (Seeley--DeWitt) expansion gives \cite{gilkey1995invariance,vassilevich2003heat}

\begin{equation}
C[g]\simeq\frac{1}{16\pi^2}\Big[f_4\Lambda^4\!\!\int\!\!\sqrt g\,\Tr\mathbf 1
+f_2\Lambda^2\!\!\int\!\!\sqrt g\,\Tr\!\Big(\tfrac R6+E\Big)
+f_0\!\!\int\!\!\sqrt g\,\Tr a_4(x)+\dots\Big],
\label{eq:HK}
\end{equation}

with moments $f_2=\int_0^\infty f\,\dd x$, $f_4=\int_0^\infty x f(x)\,\dd x$, $f_0=f(0)$.
This is the structure of the Chamseddine--Connes spectral action \cite{chamseddine1997spectral} and is related to
Sakharov's induced gravity \cite{sakharov1967vacuum}. Here the cutoff $f$ is a free choice.

\subsection{Operators}

For a Dirac operator, $D^2=\nabla^*\nabla+R/4$, so $E=-R/4$ and the curvature term in
\eqref{eq:HK} is $\tfrac R6-\tfrac R4=-\tfrac R{12}$ per component. For a minimally coupled
scalar, $E=0$ and the coefficient is $+\tfrac R6$. We write $\xi$ for the coefficient in
$L=-\nabla^2+\xi R$; $\xi=\tfrac14$ mimics the Dirac $a_0$ and $a_2$ terms (but not the spin-curvature
terms of $a_4$). Spinor dimension in 4D is $4$.

\subsection{Cost to action}

We write the Euclidean weight as $e^{-\kappa C}$ and set $\kappa=\ln 2$ in numbers, following
$P=2^{-C}$. \flag{free parameter} Matching the curvature term to the Euclidean
Einstein--Hilbert action $S_E\supset-\tfrac1{16\pi G}\!\int\! R$ gives (with $\Tr\mathbf 1=4$)

\begin{equation}
G=\frac{3\pi}{\kappa f_2\Lambda^{2}}\quad(\text{up to the stated field content}),
\label{eq:G}
\end{equation}

so $G$ is not an independent input but is tied to $\kappa f_2$.
Matching the gauge-field term from the Dirac operator in a $U(1)$ background gives
$1/e^2=\kappa f_0/(6\pi^2)$ (Section~\ref{sec:flux}).

\subsection{Assumptions register}\label{sec:assumptions}
\begin{center}
\small
\begin{tabular}{@{}l p{11.5cm}@{}}
\toprule
A1 & Smooth cutoff profile $f$; its shape and the scale $\Lambda$ are free. Hard cutoffs fail (Section~\ref{sec:sphere2}).\\
A2 & Cost-to-action factor $\kappa$ (taken as $\ln 2$ in numbers); $G\propto1/\kappa$.\\
A3 & Modes are spinor-type (this fixes the sign of the Einstein--Hilbert term). Field content (multiplicity 4) is assumed.\\
A4 & Euclidean signature, and the whole 4D history as the counted object.\\
A5 & A hand-chosen volume counterterm ($\epsilon$, or $m$) cancels the leading $\Lambda^4$ volume term. Not derived.\\
A6 & Flux sector $n$ is a fixed, quantized background; no mechanism selects it.\\
A7 & Matter: the linear cost $\omega/\Delta\omega$ acts only on \emph{occupied} modes, with $\beta=\kappa/\Delta\omega$ (Section~\ref{sec:static}).\\
A8 & In the planar test, quanta are dust-like and couple as $N\rho$ (Section~\ref{sec:planar}); not simulated.\\
A9 & Dimension $D=4$ is an input.\\
A10 & The $\xi=\tfrac14$ scalar operator stands in for Dirac in the planar test ($a_4$ spin terms not matched).\\
\bottomrule
\end{tabular}
\end{center}



\section{Tests}
\label{sec:tests}

All tests use numpy/scipy; each script prints its own tables. Scripts are in the repository
folder referenced by \texttt{\textbackslash repo} (edit the macro at the top of this file).

\begin{center}
\small
\begin{tabular}{@{}ll p{8cm}@{}}
\toprule
Section & Script & Checks\\
\midrule
\ref{sec:sphere2} & \testfile{spectral_weyl_check.py} & Weyl law (torus, inconclusive); curvature term on icosphere \\
\ref{sec:sphere2} & \testfile{sphere_cost_scan.py} & $S^2$, $C(R)$ for several $f$ \\
\ref{sec:s4} & \testfile{s4_cost_scan.py} & $S^4$, scalar vs Dirac sign of $R$-term\\
\ref{sec:s2s2} & \testfile{s2xs2_cost_scan.py} & $S^2\times S^2$, anisotropy, stationary point\\
\ref{sec:flux} & \testfile{flux_s2xs2.py} & Flux source, Einstein--Maxwell correspondence\\
\ref{sec:maxwell} & \testfile{maxwell_from_spectrum.py} & Variation of $A$; Maxwell term; gauge invariance\\
\ref{sec:static} & \testfile{two_scenarios_static_universe.py} & $S^3\times S^1$ static universe, two readings of the linear cost\\
\ref{sec:planar} & \testfile{poisson_inhomogeneous.py} & Planar inhomogeneous test, Poisson limit\\
\ref{sec:planar} & \testfile{poisson_eps_scaling.py} & Amplitude extrapolation of the planar Hessian\\
\bottomrule
\end{tabular}
\end{center}


\subsection{Two-dimensional warm-up: \texorpdfstring{$S^2$}{S2} and a discretized sphere}
\label{sec:sphere2}

\textbf{Scripts:} \testfile{spectral_weyl_check.py}, \testfile{sphere_cost_scan.py}.


In $d=2$ the curvature term $\int R$ is topological, so these runs test the machinery, not the dynamics.
For smooth $f$ the prediction is $C(R)\simeq f_2\,(\Lambda R)^2+\tfrac13 f(0)$ on a round sphere of radius $R$.
With exact spectra ($\lambda_l=l(l+1)/R^2$) and $R\in[4,14]$:

\begin{center}
\small
\begin{tabular}{@{}lcccc@{}}
\toprule
$f$ & fit $a$ & predicted $f_2$ & fit $b$ & predicted $f(0)/3$\\
\midrule
$e^{-x}$ & 1.0000 & 1.0000 & 0.336 & 0.333\\
$xe^{-x}$ & 1.0000 & 1.0000 & $-0.003$ & 0\\
$x\,\theta(1-x)$ (linear cost, hard cutoff) & 0.516 & 0.5 & 0.60 & 0\\
$x^2\theta(1-x)$ (control) & 0.349 & 0.333 & 0.61 & 0\\
$\theta(1-x)$ (hard step) & 1.016 & 1.0 & 0.92 & 0.333\\
\bottomrule
\end{tabular}
\end{center}

Smooth cutoffs match with residual $\sim\!2\times10^{-3}$. Hard cutoffs give residuals $\sim10$
from modes crossing the cutoff (Weyl fluctuations) and distorted fits.
\textbf{Consequence:} the linear cost with a sharp cutoff cannot be used to read off geometry; a smooth
cutoff is required.

On an icosphere with the cotangent Laplacian and lumped mass matrix (2562 nodes), $K(t)-1/t$ tracks
the exact sphere value ($0.362$ vs.\ $0.347$ at $t=0.2$; $0.365$ vs.\ $0.362$ at $t=0.4$), the
lowest eigenvalues are $2.000$ and $5.99$ (exact $2$, $6$), and the error at $t=0.05$ falls from
$1.28$ to $0.28$ between 642 and 2562 nodes. A discrete spectrum therefore carries the curvature
coefficient.

\textbf{Inconclusive:} the Weyl counting test on a $40\times40$ torus gave ratios $0.93$--$1.32$, and the random-graph
control gave $1.2$--$1.9$. Staircase noise from degeneracies dominates, so this test does not
discriminate geometry from no geometry at this size and should not be cited.


\subsection{\texorpdfstring{$S^4$}{S4}: scalar versus Dirac}
\label{sec:s4}
\textbf{Script:} \testfile{s4_cost_scan.py}.

On the round $S^4$ (exact scalar spectrum $l(l+3)/R^2$; Dirac $D^2$ eigenvalues $(n+2)^2/R^2$) the prediction from \eqref{eq:HK} is
\begin{align}
\text{scalar:}&\quad C\simeq\tfrac16 f_4R^4+\tfrac13 f_2R^2+f_0\tfrac{29}{90},\\
\text{Dirac:}&\quad C\simeq\tfrac23 f_4R^4-\tfrac23 f_2R^2+f_0\,a_4^{D}.
\end{align}
For $f=e^{-x}$ ($f_2=f_4=1$), fitted vs.\ predicted coefficients:
\begin{center}
\small
\begin{tabular}{@{}lcccc@{}}
\toprule
 & $R^4$ fit & $R^4$ pred. & $R^2$ fit & $R^2$ pred.\\
\midrule
scalar & 0.1667 & 0.1667 & $+0.333$ & $+0.333$\\
Dirac & 0.6667 & 0.6667 & $-0.667$ & $-0.667$\\
\bottomrule
\end{tabular}
\end{center}
Residuals are $10^{-3}$ to $10^{-5}$. \textbf{The $R^2$ term is the Einstein--Hilbert term; its sign is operator dependent:} the scalar spectrum gives the wrong
sign for Euclidean gravity ($S_E\supset-\int R/16\pi G$), the spinor spectrum the right one.
The control $f=x^2e^{-x}$ obeys the same expansion with different coefficients, so the check
validates the machinery, not the choice of $f$.

\textbf{No usable extremum.} For Dirac the asymptotic formula has a stationary point at
$R^2\Lambda^2=f_2/(2f_4)=0.5$, outside the regime where the expansion holds; the exact $C(R)$ tends to $\approx0$ as
$R\to0$. In this measure small universes are the cheapest. The volume term is $O(\Lambda^4)$ and the
curvature term $O(\Lambda^2)$, so any saddle sits at Planck size unless the volume term is cancelled (A5).


\subsection{\texorpdfstring{$S^2\times S^2$}{S2xS2}: anisotropy and a stationary point}
\label{sec:s2s2}
\textbf{Script:} \testfile{s2xs2_cost_scan.py}.

With $x=r_1^2$, $y=r_2^2$, the prediction is
\begin{align}
\text{scalar:}&\quad C\simeq f_4xy+\tfrac{f_2}{3}(x+y)+a_4,\\
\text{Dirac:}&\quad C\simeq 4f_4xy-\tfrac{2f_2}{3}(x+y)+a_4.
\end{align}

For $f=xe^{-x}$ ($f_2=1,f_4=2$, $f(0)=0$ so $a_4$ drops) the fits over $r_{1,2}\in[3,9]$ give
$a=2.0000$, $b=+0.333$ (scalar) and $a=8.0000$, $b=-0.6665$ (Dirac), matching $f_4$, $f_2/3$, $4f_4$, $-2f_2/3$.
With $C=a\,xy+b(x+y)$ the stationary conditions require $x=y=-b/a$, which needs $b<0$: only the Dirac sign has a
stationary point. The pure cost puts it at Planck size, so we subtract a hand-chosen volume term
$\mu\,xy$ with $\mu=a(1-\epsilon)$ \flag{A5, free parameter}. Then, for $\epsilon=0.01$ and $0.004$, the exact stationary points
are $x=y=8.33$ and $20.83$ (analytic $-b/(a\epsilon)$), with Hessian eigenvalues $\mp 0.08$ and $\mp 0.032$.
There is one negative direction (volume-preserving anisotropy), a feature I believe is
shared by Euclidean Einstein gravity, but I have not checked this against the literature.
The scalar case has no stationary point (the solver did not converge).
The stationary point at $r_1=r_2$ is guaranteed by symmetry; the informative parts are the sign
requirement and the Hessian.


\subsection{A flux source: Einstein--Maxwell correspondence}
\label{sec:flux}
\textbf{Script:} \testfile{flux_s2xs2.py}.

Put a $U(1)$ magnetic flux $n$ through $S^2(r_2)$ (field strength $B=n/2y$) and use the exact charged Dirac spectrum
(on $S^2$: $k(k+|n|)/y$, multiplicity $2(2k+|n|)$, plus $|n|$ zero modes; this standard
monopole result is validated by the fit below). Evaluating the heat-kernel $a_4$ coefficient (standard Gilkey formula; my own evaluation)
for the Dirac operator on $S^2\times S^2$ gives
\begin{equation}
C\simeq4f_4xy-\tfrac23f_2(x+y)+f_0\Big[\tfrac{40-24(x/y+y/x)}{360}+\tfrac{n^2x}{3y}\Big].
\label{eq:Cflux}
\end{equation}
The gauge-field term in the $a_4$ density is $\tfrac1{16\pi^2}\tfrac23 f_{\mu\nu}f^{\mu\nu}$, so
$\Delta C=\tfrac{f_0}{24\pi^2}\int f_{\mu\nu}f^{\mu\nu}$. For $f=e^{-x}$, the difference between exact and
\eqref{eq:Cflux} is $\sim10^{-3}$ at $n=0$ and at most $0.04$ at $n=3$.

\textbf{Stationary points} of $C-\mu\,xy$ with $\epsilon=0.01$ (A5, A6):

\begin{center}
\small
\begin{tabular}{@{}cccc@{}}
\toprule
$n$ & analytic $(x,y)$ & exact $(x,y)$ & $x/y$\\
\midrule
0 & (16.67, 16.67) & (16.67, 16.67) & 1.000\\
1 & (17.22, 16.15) & (17.21, 16.16) & 1.065\\
2 & (19.89, 14.34) & (19.85, 14.37) & 1.381\\
\bottomrule
\end{tabular}
\end{center}

Flux on $S^2(r_2)$ shrinks it and enlarges the other sphere, the direction required by Einstein--Maxwell.
The Hessian keeps one negative eigenvalue. Comparing with the Euclidean Einstein--Maxwell equations
(which I wrote myself and which deserve an independent check) gives
$1/x=1/y-2\pi G n^2/(e^2y^2)$, matching the cost iff

\begin{equation}
\frac{2\pi G}{e^2}=\frac{f_0}{f_2}.
\label{eq:Ge}
\end{equation}

An independent action-level match gives $G=3\pi/(\kappa f_2)$ and $1/e^2=\kappa f_0/(6\pi^2)$; the ratio
reproduces \eqref{eq:Ge} and $\kappa$ drops out.

\textbf{Control ($f=xe^{-x}$, $f(0)=0$).} The Maxwell term sits at $a_4$ and multiplies $f(0)$, so no leading-order
flux effect is expected. At $n=2$ the exact result is $x/y=0.988$: small, in the opposite direction, attributed (but not shown) to a
subleading term. \textbf{Gravity and matter couple through different moments of $f$:} curvature and volume through
$f_2,f_4$; the gauge term through $f(0)$.

\textbf{Limits.} The flux $n$ is a fixed quantized background; the gauge field was not varied here;
$n$ is not selected by anything; at fixed geometry, $\Delta C=f_0n^2x/(3y)$ would suppress flux
sectors by $2^{-n^2x/3y}$ (assuming $\kappa=\ln2$), a mild effect not tested jointly with geometry readjustment.


\subsection{Varying the gauge field: Maxwell from the spectrum}
\label{sec:maxwell}
\textbf{Script:} \testfile{maxwell_from_spectrum.py}.

On flat $T^2(L)\times T^2(L)$ ($L=20$, no curvature) a charge-1 field $A_2=a\cos(kx_1)$ (transverse) enters the Dirac operator.
With $f=e^{-x}$ the cost factorizes and the exact spectrum is obtained by Fourier-basis
diagonalization (sanity check against the $A=0$ trace: agreement to 8 digits). The prediction is
$\Delta C=f_0L^4a^2k^2/(24\pi^2)$.

\begin{center}
\small
\begin{tabular}{@{}ccc@{}}
\toprule
$k$ & $\Delta C/\mathrm{pred}$ ($f=e^{-x}$) & $\Delta C/\mathrm{pred}$ (control $f=xe^{-x}$)\\
\midrule
0.314 & 0.9805 & 0.019\\
0.628 & 0.9243 & 0.073\\
0.942 & 0.838 & 0.147\\
\bottomrule
\end{tabular}
\end{center}

Doubling the amplitude $a$ quadruples $\Delta C$ to four digits. A pure-gauge field $A_1=a\cos(kx_1)$ gives $\Delta C\approx10^{-12}$.
The response is $\propto k^2$ for transverse fields and zero for longitudinal ones: the Maxwell structure
$k^2\delta^{\mu\nu}-k^\mu k^\nu$. The quadratic form is positive for $k\neq0$, so the only stationary point is $a=0$ and the field equation is $\partial_\mu f^{\mu\nu}=0$.
The ratio falls roughly as $1-0.2\,k^2/\Lambda^2$ (a higher-derivative correction not derived here); the control vanishes as $k\to0$
as expected for $f(0)=0$.
\textbf{Limits:} only the quadratic response about $A=0$ on flat tori; no charged sources, no non-abelian fields, no nonlinear response.


\subsection{Static universe: two readings of the per-mode cost}
\label{sec:static}

\textbf{Script:} \testfile{two_scenarios_static_universe.py}.

\paragraph{The mismatch.} The cost \eqref{eq:Cs} sums over the active modes of a state, whereas the spectral cost sums over all
eigenmodes of the geometry's operator. We test two bridges on $S^3(r)\times S^1(\beta)$ with $\beta\gg1/\Lambda$,
Dirac modes $\omega_n=(n+\tfrac32)/r$:
\begin{description}
\item[Scenario 1:] full spectral geometry term with smooth $f$, plus the linear cost $N\omega_0/\Delta\omega$ on $N$ \emph{occupied} quanta
of the lowest mode, with the identification $\beta=\kappa/\Delta\omega$ \flag{A7, assumption}. The identification makes the matter
cost proportional to $\beta$, so varying $\beta$ gives the Hamiltonian constraint with the matter energy in it.
\item[Scenario 2:] the linear cost on \emph{all} modes with no separate geometry term (vacuum cost $E(r)=\sum_n 2(n+1)(n+2)\,\omega_ne^{-\omega_n^2}$).
\end{description}

\paragraph{Scenario 1.} The exact geometry term is $e_1(r)=\tfrac12r^3-\tfrac14r$ (no $1/r$ or $1/r^3$ corrections at the 5-digit level),
matching the heat-kernel coefficients. Imposing both stationarity conditions (with the counterterm $m$ solved for) gives
the Einstein-static-universe radius $r^2=12N/(\kappa f_2)$ with $G$ from \eqref{eq:G}:

\begin{center}
\small
\begin{tabular}{@{}ccc@{}}
\toprule
$N$ & $r$ (exact) & $r^2/r^2_{\mathrm{Einstein}}$\\
\midrule
5 & 9.304 & 1.0000\\
10 & 13.158 & 1.0000\\
20 & 18.608 & 1.0000\\
40 & 26.315 & 1.0000\\
\bottomrule
\end{tabular}
\end{center}

\textbf{Note:} The ratio is exactly 1 because $e_1$ has no corrections, so this is an algebraic check: the two stationarity conditions
reproduce the Hamiltonian constraint and the pressure equation, with the same $G$ in the curvature term and the source coupling. It is not independent evidence
for the linear cost. A matter cost scaling as $r^{-s}$ passes for $s=0$ (dust-like, $r=129.843$), $s=1$ (radiation-like, $13.158$) and
$s=2$ ($6.004$), each equal to the Einstein-static radius, at $N=10$. So the test does not single out linear-in-$\omega$: whatever cost is
proportional to $\beta$ gravitates as energy. If $\Delta\omega$ were independent of $\beta$, the $\beta$-equation would contain no matter
(argued algebraically, not run).

\paragraph{Scenario 2.} The vacuum cost fits $E(r)=r^3-\tfrac14r+c/r$, with the same Einstein--Hilbert-like coefficient $-\tfrac14$ and
a Casimir-like $c=-0.0354$, matching $-17/480$ once the third Euler--Maclaurin term is included. With the $\beta$-equation imposed,
the $r$-condition is $F(r)=rE'-3E=-2Br-4c/r$ with $B=-\tfrac14$ and $c<0$, i.e.\ both terms positive. Numerically $F=0.82,\,2.53,\,5.01,\,10.0$ at $r=1,5,10,20$ and there is no root on $[0.2,20]$:
\textbf{no static solution.} The vacuum term adds to the curvature term and does not oppose it. The sign of $c$
depends on the field content (Dirac counting assumed here), so I do not claim scenario 2 is excluded for every field content. Scenario 2 also has no tunable matter and, with $f(0)=0$, no Maxwell term at leading order.

\paragraph{Conclusion.} Scenario 1 works with three requirements: a distinction between occupied and unoccupied modes, the identification
$\beta=\kappa/\Delta\omega$, and a separate smooth $f$ with $f(0)\neq0$ for the gauge sector.


\subsection{Planar inhomogeneous test: the Poisson limit}
\label{sec:planar}
\textbf{Scripts:} \testfile{poisson_inhomogeneous.py}, \testfile{poisson_eps_scaling.py}.

\paragraph{Setup.} Metric $ds^2=N(z)^2\dd\tau^2+e^{2\sigma(z)}(\dd x^2+\dd y^2+\dd z^2)$, long $\tau$-circle, large transverse area, $z$ periodic with $L_z=40$.
The operator is $L=-\nabla_4^2+\xi R_4$ with $\xi=\tfrac14$ and multiplicity $4$ (A10), $f=e^{-x}$. With $N=1+\delta\cos kz$ and $\sigma=\epsilon\cos kz$, the heat-kernel prediction
is $C/(\beta A_\perp)\simeq\tfrac1{4\pi^2}\!\int\!\dd z\,Ne^{3\sigma}[f_4+f_2(\tfrac16-\xi)R_3]$ (the total derivative in $R_4$ drops). After
tuning away the volume term (A5), the Hessian at amplitude $\cos kz$ is
\begin{equation}
C_{\epsilon\delta}=C_{\epsilon\epsilon}=\frac{L_z}{8\pi^2}\,4f_2\big(\tfrac16-\xi\big)k^2,\qquad C_{\delta\delta}=0 .
\end{equation}
The spectrum is computed with a Fourier-pseudospectral discretization in $z$ and Gauss quadrature over $\omega$ and $k_\perp$.
Sanity check $C_0=L_z/4\pi^2$ agrees to 14 digits (after fixing the Nyquist artifact, Section~\ref{sec:failures}).

\paragraph{Results ($\xi=\tfrac14$), amplitude-extrapolated.}
\begin{center}
\small
\begin{tabular}{@{}cccc@{}}
\toprule
$k$ & $C_{\epsilon\delta}$ (units $L_z/8\pi^2$) & predicted $-k^2/3$ & $C_{\epsilon\epsilon}/C_{\epsilon\delta}$\\
\midrule
0.314 & $-0.0326$ & $-0.0329$ & 0.970\\
0.471 & $-0.0724$ & $-0.0740$ & 0.933\\
0.628 & $-0.1265$ & $-0.1316$ & 0.880\\
\bottomrule
\end{tabular}
\end{center}
The $k^2$ slope extrapolates to $f_{2,\mathrm{eff}}=1.00$, so $G_{\mathrm{eff}}=3\pi/(\kappa f_2)\approx13.6$ (cutoff units, $\kappa=\ln2$): the same $G$ as in the static test. With matter coupled as
$N\rho$ (A8), the constraint gives $\sigma=4\pi G\rho_a/k^2$; with $\Phi=N-1$ this is $\nabla^2\Phi=4\pi G\rho$, an attractive potential.
The $\sigma$-equation requires $C_{\epsilon\epsilon}=C_{\epsilon\delta}$ (Newtonian gauge $N-1=-\sigma$); the ratio tends to 1 as $k\to0$. Also $C_{\delta\delta}=0.0487\,k^4$:
linear in the lapse at leading order, with an $a_4$-type $k^4$ correction.
Deviations grow like $k^2/\Lambda^2$ (about $0.1k^2$ for $C_{\epsilon\delta}$, $0.3k^2$ for the ratio), so Poisson holds only for wavelengths well above $1/\Lambda$.

\paragraph{Controls.} $\xi=0$ (minimal scalar): $C_{\epsilon\delta}>0$ (about $+\tfrac23k^2$), i.e.\ repulsive gravity.
$\xi=\tfrac16$ (conformal): no systematic $k^2$ term, i.e.\ no gravitational coupling. The check discriminates, and the spinor-type
curvature coupling sets the sign.

\paragraph{Limits.} Planar symmetry and weak field only; matter not simulated (assumed to couple as $N\rho_{\mathrm{loc}}$,
dust-like); no point source with $1/r$ falloff; no nonlinear response; no Lorentzian dynamics.


\section{Summary of relations obtained}
\label{sec:summary}
Under Assumptions A1--A10, with $\Lambda=1$:

\begin{align}
G&=\frac{3\pi}{\kappa f_2}, & \frac1{e^2}&=\frac{\kappa f_0}{6\pi^2}, & \frac{2\pi G}{e^2}&=\frac{f_0}{f_2},\\
\Delta C_{\mathrm{gauge}}&=\frac{f_0}{24\pi^2}\!\int\! f_{\mu\nu}f^{\mu\nu}, &
r^2_{\mathrm{static}}&=\frac{12N}{\kappa f_2}, & \nabla^2\Phi&=4\pi G\rho.
\end{align}

Gravity (curvature and volume) is carried by the energy-weighted moments $f_2,f_4$ and the gauge coupling by $f(0)$; the
ratio $G/e^2$ is $\kappa$-independent and fixed by the cutoff profile.


\section{Failures, corrections and weak points}
\label{sec:failures}
This is a complete log of what did not work or was corrected during the work.
\begin{enumerate}
\item \textbf{Hard cutoffs fail} (Section~\ref{sec:sphere2}): sawtooth from mode crossings, wrong constant term; the linear cost needs a smooth cutoff.
\item \textbf{Weyl counting test inconclusive} on a $40\times40$ torus; the random-graph control did not separate (ratios $1.2$--$1.9$ vs.\ $0.93$--$1.32$).
\item \textbf{No usable extremum} on $S^2$/$S^4$ without the hand-tuned volume counterterm: the exact $C(R)\to0$ as $R\to0$; the asymptotic stationary
point at $R^2\Lambda^2=0.5$ lies outside the expansion's validity. The cosmological-constant problem is not addressed; any saddle is Planck sized unless A5 holds.
I did not separate the non-monotone sawtooth of hard-cutoff $C(R)$ from a real minimum.
\item \textbf{Scalar operator gives the wrong sign} of the curvature term (no stationary geometry on $S^2\times S^2$; repulsive in the planar test).
\item \textbf{Control $f=xe^{-x}$ in the flux test} gave $x/y=0.988$ at $n=2$ (small, opposite direction) rather than exactly $1$; attributed, not derived, to a subleading term.
\item \textbf{Static-universe agreement is algebraic}: $e_1(r)$ has no corrections, and the controls $s=0,1,2$ also pass, so the test does not single out linear-in-$\omega$.
\item \textbf{Scenario 2 has no static solution} (for Dirac counting).
\item \textbf{Higher-derivative corrections} in the Maxwell and Poisson tests ($\sim0.2k^2$, $0.1k^2$, $0.3k^2$) were measured, not derived.
\item \textbf{Euclidean instability:} the stationary points have a negative Hessian direction; I believe this matches Euclidean Einstein gravity but did not verify.
\end{enumerate}
\textbf{Errors of mine that were caught and fixed} (included for the record; the numbers above are post-correction):
\begin{itemize}
\item In the planar test, an even Fourier grid left a spurious Nyquist zero mode, giving a 9\% error in the sanity check; fixed with an odd grid, after which the check agrees to 14 digits. The first run's numbers were discarded.
\item The first run of the static-universe script crashed in the controls section (root bracket too narrow), so the scenario-2 search did not run; fixed and rerun.
\item My first Euler--Maclaurin estimate of the Casimir-like coefficient ($-1/48$) omitted a term; the fit gives $-0.0354=-17/480$. I also wrote the scenario-2 stationarity condition with a wrong sign for the $B$ term at first; the corrected form shows no root.
\end{itemize}
\textbf{Not independently verified:} the Gilkey $a_4$ evaluation (self-computed, validated only by the numerical fits above), the Euclidean Einstein--Maxwell equations
used in Section~\ref{sec:flux}, the monopole-Dirac spectrum (validated through the fit), and the claim that the sign structure matches the Chamseddine--Connes result (recalled from the literature).


\section{What has not been shown}
\label{sec:open}

\begin{itemize}
\item A derivation of the cutoff $f$, $\Lambda$, $\kappa$ or the cosmological cancellation $\epsilon$.
\item Lorentzian signature and causality; the Euclidean cost over the whole history is assumed.
\item The origin of $D=4$.
\item The occupied-mode prescription and $\beta=\kappa/\Delta\omega$ from the simulation's observer filter; $N$ is a free input.
\item Full Einstein equations with dynamical matter; non-planar sources; nonlinear response; non-abelian gauge fields; charged sources; selection of the flux sector.
\item Any test of the lognormal expansion picture.
\end{itemize}


\section{Proposed stress tests}
\label{sec:stress}
\begin{enumerate}
\item \textbf{Cutoff dependence.} Repeat all tests for several smooth $f$ (including ones with $f(0)=0$, $f(0)\neq0$, different $f_2/f_4$); quantify which conclusions are cutoff independent.
\item \textbf{Independent verification of formulas.} Re-derive the $a_4$ gauge and curvature coefficients symbolically and against Vassilevich's tables; check the Euclidean Einstein--Maxwell equations independently.
\item \textbf{Other geometries.} Non-Einstein metrics, $S^2\times S^2$ with unequal gauge flux, $\mathbb{CP}^2$, squashed spheres; spherically symmetric point sources ($1/r$).
\item \textbf{Dynamical gauge sector.} Sum over flux sectors with the exact weights; vary $A$ nonperturbatively; non-abelian extension.
\item \textbf{Numerics robustness.} Grid, quadrature and basis-size convergence; confirm that $O(\epsilon^2)$ and $O(k^2)$ corrections scale as claimed.
\item \textbf{Higher-order terms.} Compute the $a_6$ corrections and compare with the measured $0.1$--$0.3\,k^2$ deviations.
\item \textbf{Occupied versus all modes on real states.} Apply the cost literally to the simulation's sampled wavefunctions and see whether the $\beta=\kappa/\Delta\omega$ identification emerges or must be imposed.
\item \textbf{Observer filter and the cosmological term.} Determine whether the observer-filter selection can supply $\epsilon$ and the lognormal time dependence.
\item \textbf{Dimension dependence.} Repeat with $D=3,5,6$ using a documented normalization convention to see whether $D=4$ is singled out.
\item \textbf{Lorentzian signature and stability.} Study the negative Hessian modes and the conformal-factor problem; attempt a Lorentzian treatment.
\item \textbf{Adversarial controls.} For every ``conservation''-type or consistency check, evaluate it on a deliberately wrong alternative and flag tautologies, as done here for $\xi$, $f(0)$, $s$ and the matter-cost scaling.
\end{enumerate}


\section{Reproducibility}

All scripts use \texttt{numpy} and \texttt{scipy} (Python~3). Each is run as \texttt{python <script>} and prints its own tables.
Runtimes are seconds to a few minutes. Script-to-section mapping is given in Section~\ref{sec:tests}; edit the macro
\texttt{\textbackslash repo} at the top of this file to point to your repository.

\section*{Acknowledgments}

The numerical exploration and drafting of this note were assisted by Claude (Anthropic). The scientific direction,
assumptions and framework are the author's. The results have not been independently reviewed.

\end{document}
